\subsection{Evaluation Metrics}
\label{sec:metrics}

Let $y_i$ be a reference GHI value, $\widehat{y}_i$ its post-processed
forecast, $e_i=y_i-\widehat{y}_i$, $\bar{y}$ the mean reference value in the
same evaluation vector, and $n$ its number of origin--lead forecast pairs.
A target timestamp can therefore occur more than once when it is predicted
from different origins. Each metric is
computed first for one site, one model estimate, one partition, one declared
horizon, and one horizon scope.

\subsubsection{Mean absolute error}

\begin{equation}
 \mathrm{MAE}
 =
 \frac{1}{n}\sum_{i=1}^{n}\left|e_i\right|.
 \label{eq:mae}
\end{equation}
MAE is the primary error measure and remains in W/m$^2$. Its linear penalty
makes the magnitude directly interpretable and less dominated by individual
large deviations than a squared-error score \cite{hyndman2006}. Lower values
are better.

\subsubsection{Root mean squared error}

\begin{equation}
 \mathrm{RMSE}
 =
 \sqrt{\frac{1}{n}\sum_{i=1}^{n}e_i^2}.
 \label{eq:rmse}
\end{equation}
RMSE is also expressed in W/m$^2$, but squaring gives larger deviations more
weight. Reporting it alongside MAE reveals whether occasional large errors
change the model ordering, a standard concern in deterministic solar-forecast
verification \cite{yang2020verification}. Lower values are better.

\subsubsection{Normalized root mean squared error}

\begin{equation}
 \mathrm{nRMSE}
 =
 \frac{\mathrm{RMSE}}{\bar{y}},
 \qquad
 \mathrm{nRMSE}_{\%}=100\,\mathrm{nRMSE}.
 \label{eq:nrmse}
\end{equation}
The machine-readable artifacts store the dimensionless ratio; tables that use
percent explicitly multiply it by 100. Normalization supports comparison among
sites with different mean irradiance levels, but it inherits the instability
of a small denominator \cite{voyant2022}. The value is therefore left undefined
when $\bar{y}$ is numerically zero. Lower values are better.

\subsubsection{Coefficient of determination}

\begin{equation}
 R^2
 =
 1-
 \frac{\sum_{i=1}^{n}e_i^2}
      {\sum_{i=1}^{n}(y_i-\bar{y})^2}.
 \label{eq:r2}
\end{equation}
$R^2$ measures variation tracked relative to a constant predictor built from
the evaluation mean \cite{yang2020verification}. It is undefined for fewer
than two targets or a constant reference vector and can be negative when the
forecast is worse than that constant reference. Higher values are better;
$R^2$ does not replace the scale-dependent errors.

\subsubsection{Horizon scopes and spatial aggregation}

For a declared horizon $H$, the cumulative scope concatenates leads
$1,\ldots,H$ over all valid origins, whereas the exact-lead scope retains only
lead $H$. The two scopes are identical at $H=1$. This distinction is applied at
24, 48, and 72 hours; 1, 3, 7, 14, and 30 days; and, for the exploratory
multi-month task, 1, 3, and 6 months. A horizon label therefore always denotes
both a numerical lead and an explicitly named scope. The main rankings use
the cumulative scope.

Hourly origins are fixed at local midnight, so exact leads 24, 48, and 72
all select the 23:00 interval on their respective target days. In this test
set these nighttime targets and post-processed forecasts are zero: exact-lead
MAE and RMSE are zero for every model, while nRMSE and $R^2$ are undefined.
They are retained as audit outputs and do not measure daytime forecast skill.

Within each site, cumulative metrics weight origin--lead pairs equally,
including repeated forecasts for a shared target timestamp. They therefore
evaluate the issued forecast trajectories rather than one unique forecast
per calendar time. Target periods inside an overlapping window can receive
more contributions than periods near the test boundaries.

Let $\mathcal{L}_{M}$ be the set of locations with a finite site-level value
of metric $M$, and let $L_M=|\mathcal{L}_{M}|$. The macro result is
\begin{equation}
 M_{\mathrm{macro}}
 =
 \frac{1}{L_M}\sum_{\ell\in\mathcal{L}_{M}}M_{\ell},
 \qquad L_M>0.
 \label{eq:macro_metric}
\end{equation}
Undefined site-level values are excluded from the corresponding mean; if
none are finite, the macro value remains undefined. All ten sites contribute
to the main cumulative results. Equal site weighting prevents a site with more
valid forecast pairs from dominating the geographical summary.
Macro-MAE is the predeclared primary
metric at every resolution; MAE, RMSE, nRMSE, and $R^2$ are nevertheless
retained together rather than selectively reported.

For daily and monthly learned models, seed-level metrics describe variability
across five fitted instances on the same test data. Their standard deviation
measures sensitivity to the registered seeds; it is neither a forecast
prediction interval nor a confidence interval for generalization error.
The TimesNet--DilatedRNN paired comparison uses the post-processed five-seed
ensemble forecast from Eq.~\eqref{eq:ensemble_forecast}. The hourly extension compares the single
registered seed 42 for both models. For each site, metric, horizon, and scope,
the pipeline stores
\begin{equation}
 \Delta_{\ell,M}
 =
 M_{\ell,\mathrm{TimesNet}}-M_{\ell,\mathrm{DilatedRNN}}.
 \label{eq:paired_difference}
\end{equation}
A negative $\Delta_{\ell,M}$ is a TimesNet win for the three error metrics; a
positive value is a win for $R^2$. Mean and median differences and the number
of site wins are descriptive. Overlapping multi-step windows are not treated
as independent observations for an i.i.d. significance test.
